\documentclass{article}
\usepackage[T1]{fontenc}
\usepackage{lmodern}
\usepackage{graphicx}
\usepackage{amsmath,amssymb}
\usepackage[a4paper,margin=2cm]{geometry}
\usepackage[absolute,overlay]{textpos}
\setlength{\TPHorizModule}{1mm}
\setlength{\TPVertModule}{1mm}
\usepackage[indonesian]{babel}
\usepackage{hyperref}
\setlength{\emergencystretch}{2em}
% unit: Halaman 28

\begin{document}

% === Halaman 28 (llm) ===

\section*{Bentuk umum \textit{cipher} abjad-majemuk}

\begin{itemize}
    \item Kunci:
    \[ K = k_1 k_2 \dots k_m \quad \text{(ket: } m = \text{ panjang kunci)} \]
    Karena $m$ lebih pendek daripada panjang plainteks, maka kunci diulang secara periodik
    
    \item Plainteks:
    \[ P = p_1 p_2 \dots p_m p_{m+1} \dots p_{2m} \dots \]
    
    \item Cipherteks:
    \[ C = E_k(P) = f_{k_1}(p_1) f_{k_2}(p_2) \dots f_{k_m}(p_m) f_{k_1}(p_{m+1}) f_{k_2}(p_{m+2}) \dots f_{k_m}(p_{2m}) \dots \]
    $f(.)$ adalah fungsi enkripsi pada \textit{cipher} abjad-tunggal

    \item Untuk $m = 1$, \textit{cipher}-nya ekivalen dengan \textit{cipher} abjad-tunggal.
\end{itemize}

\end{document}
